Pro Kilogramm muss das Eisen zuerst von $\vartheta_1 = \TaO$ auf den Schmelzpunkt $\vartheta_\text{f} = \TfFeO$ erwärmt und dann geschmolzen werden. Mit $c = \cFeO$ und $L_\text{f} = \LfFeO$ (FoTaBe S.~190) braucht das Erwärmen pro Kilogramm
\begin{align}
 \frac{Q_1}{m} &= \QaFeF = \cFe\cdot\left(\TfFeO-\TaO\right) = \QaFe
\end{align}
Der Anteil des Phasenübergangs an der gesamten Wärme ist
\begin{align}
 \frac{Q_\text{f}}{Q_1+Q_\text{f}} &= \antF \\
   &= \frac{\LfFe}{\QaFe+\LfFe} \\
   &= \ant \approx \result{\antP{0}{2}}
\end{align}
Gerechnet ist mit der spezifischen Wärme bei Zimmertemperatur. Diese nimmt bei hohen Temperaturen zu, der tatsächliche Anteil der Schmelzwärme ist deshalb etwas kleiner.
